\section{Vulnerabilities and Malware Enablers}

To better understand how malware operates, it is crucial to examine the methods and techniques used by attackers. These techniques exploit system vulnerabilities, allowing malicious code to infiltrate systems and cause harm. In this section, we explore common attack methods and analyze how malware targets computer systems and networks.

\subsection{Active Content Vulnerabilities}

\paragraph{Definition.}
Active content refers to components, primarily on web pages, that provide functionality for interacting with users. This includes any dynamic object that performs an action when a user opens a web page. Common technologies used to implement active content include ActiveX, Java, JavaScript, VBScript, scripts, browser plug-ins, and PDF files.

\paragraph{How It Works.}
Active content is considered a form of mobile code because it can execute across different computing platforms. When a user visits a website, the web browser downloads this code and executes it within the browser's context, using the user's credentials and privileges. Once downloaded, the mobile code may gain access to local system resources, including the computer's hard drive.

\paragraph{Attack Techniques.}
Attackers exploit weaknesses in active content mechanisms to execute malicious code. A common attack technique is Cross-Site Scripting (XSS), which allows attackers to inject client-side scripts into web pages viewed by users. When a user accesses such a page, the browser automatically executes the malicious script, often bypassing access control mechanisms.

\paragraph{Impact.}
Malicious active content can cause severe consequences, such as flooding the desktop with icons representing malware-infected files. These files may continue to replicate and generate additional copies of the malicious code. The overall impact depends on the sensitivity of the compromised system and the data it contains.

\paragraph{Prevention.}
\begin{itemize}
  \item Browser configuration: Configure browser security settings to block or warn users before executing scripts or dynamic content.
  \item Content filtering: Implement network-level filtering to analyze traffic, identify dynamic code components (e.g., Java or ActiveX), and disable script processing within web browsers.
  \item Pop-up blocking: Require users to enable pop-up blockers to prevent automatic pop-up windows and reduce the risk of unintentional malware execution.
\end{itemize}

\subsection{Homepage Hijacking}

\paragraph{Definition.}
Homepage hijacking is a type of attack that alters a web browser's homepage configuration so that it redirects users to a website chosen by the attacker.

\paragraph{How It Works.}
This attack resets the browser's homepage settings, typically without the user's permission. Attackers may install hidden programs that preserve the malicious configuration even when the user attempts to restore the original homepage.

\paragraph{Attack Techniques.}
Attackers commonly use the following methods to perform homepage hijacking:

\begin{itemize}
  \item Exploiting browser vulnerabilities: Taking advantage of security flaws in web browsers or active content components to modify homepage settings.
  \item Social engineering: Tricking users into performing actions that appear beneficial but actually initiate malware execution (e.g., fake buttons claiming to remove infections).
  \item Browser Helper Objects (BHOs): Installing Trojan programs in the form of BHOs that register themselves in the system startup process and execute automatically each time the system boots.
\end{itemize}

\paragraph{Impact.}
\begin{itemize}
  \item The browser repeatedly opens the attacker's website instead of the user's preferred homepage.
  \item Persistent annoyance and loss of control, as hijacking programs continuously reset the homepage until the malware is fully removed.
\end{itemize}

\paragraph{Prevention.}
\begin{itemize}
  \item Exercise caution when interacting with websites, especially when clicking buttons or links from unknown or untrusted sources.
  \item Review system startup processes and remove suspicious or unauthorized entries.
  \item Fully identify and eliminate malicious programs or Browser Helper Objects to prevent repeated homepage changes.
\end{itemize}

\subsection{Webpage Defacements}

\paragraph{Definition.}
Webpage defacement (also known as web graffiti) is the act of gaining unauthorized access to a web server and altering the content of the index page or other web pages hosted on that server.

\paragraph{How It Works.}
Attackers typically target publicly accessible web servers on the Internet. The attack process involves identifying and exploiting known security vulnerabilities to gain administrative or elevated access. Once control is obtained, the attackers replace legitimate web pages with modified or malicious versions.

\paragraph{Attack Techniques.}
\begin{itemize}
  \item Exploiting software vulnerabilities: Attackers exploit security flaws in the server's operating system or service applications (e.g., vulnerabilities in web servers) to gain unauthorized access.
  \item Installing backdoors: Some malware installs Trojan backdoors after infection, enabling attackers to remotely access the server and alter website content at any time.
  \item Using automated tools: Script kiddies and organized hacker groups often use automated vulnerability-scanning tools to identify poorly secured websites before launching attacks.
\end{itemize}

\paragraph{Impact.}
\begin{itemize}
  \item Defacement and nuisance: Webpage defacement is often considered a form of digital graffiti that disrupts normal website operations.
  \item Reputational damage: Defaced websites may display offensive messages or inappropriate content, damaging organizational credibility and user trust.
  \item Serious security risks: Attackers may embed malicious active content, such as viruses or Trojan programs, exposing visitors to further compromise.
\end{itemize}

\paragraph{Prevention.}
\begin{itemize}
  \item Software updates: Keep web server operating systems and applications up to date with the latest security patches.
  \item Disabling unnecessary services: Turn off unused services and close unneeded network ports to reduce attack surfaces.
  \item Network defense layers: Use proxy servers or bastion hosts to protect critical services and filter malicious traffic.
  \item Monitoring and testing: Conduct regular vulnerability assessments and penetration testing to identify and remediate weaknesses proactively.
\end{itemize}
